\noindent\textbf{Elección del procedimiento.} La viscosidad cinemática se define como $\nu=\mu/\rho$. Se calcula por separado en cada fluido y después se divide; comparar solo viscosidades dinámicas omite el efecto de la densidad.

\begin{enumerate}

\item Calcular la viscosidad cinemática del aire.

\[\fitmath{\mathrm{nua}=\frac{\mathrm{mua}}{\mathrm{rhoa}}}\]

\[\fitmath{\mathrm{nua}=\frac{(1.7894\times10^{-5})}{(1.225)}\approx 1.460735\times10^{-5}\,\mathrm{m^{2}/s}}\]

\item Calcular la del agua con su propia densidad.

\[\fitmath{\mathrm{nuw}=\frac{\mathrm{muw}}{\mathrm{rhow}}}\]

\[\fitmath{\mathrm{nuw}=\frac{(0.001005)}{(998.2)}\approx 1.006812\times10^{-6}\,\mathrm{m^{2}/s}}\]

\item Dividir las magnitudes de igual dimensión.

\[\fitmath{\mathrm{ratio}=\frac{\mathrm{nua}}{\mathrm{nuw}}}\]

\[\fitmath{\mathrm{ratio}=\frac{(1.460734694\times10^{-5})}{(1.006812262\times10^{-6})}\approx 14.50851\,}\]

\end{enumerate}
